\section{\label{II_3_C}Rendre les données recherchables et découvrables} 

Une fois les données validées et indexées au sein du modèle de données, elles modifient l'accès aux archives. L'indexation de séquences décrites par une nouvelle entité de description, permet à  l'institution de non plus seulement avoir accès aux archives dans leur globalité, mais à une "indexation fine" de leur contenu \footcite{bachimont2007}. Ainsi, par une recherche de mots-clés, l'utilisateur peut directement accéder aux segments d’archives concernés. La modélisation descriptive prévue par le FRBR permet de conserver le contexte du fragment, associé d'une part à l'oeuvre dont il est extrait, et d'autre part au média numérique associé. Les liens entre le fragment et les entités liées sont toujours requis, empêchant ainsi un risque de désolidarisation du contexte. 
Les séquences créées par le processus automatique pourront ainsi être requêtées par l'utilisateur via les outils de recherche offerts par le logiciel développé par Skinsoft. En effet, le logiciel propose les modes de recherche suivants : recherche simple, recherche avancée, recherche plein texte ciblée, recherche par image. De plus, une recherche en langage naturel, basée sur le \gls{nlp} et sur une génération augmentée de récupération (\gls{rag}), imaginée au cours de notre stage, est en cours de développement au sein du logiciel.  
Ce mode de recherche permet à l'utilisateur de formuler des requêtes complexes en langage naturel, afin d'améliorer l'expérience de navigation au sein de l'interface. L'utilisateur, via ce mode de recherche, n'a plus besoin d'employer les mots-clés exacts pour obtenir certains résultats. Par exemple, il peut formuler la requête suivante : "toutes les archives sur des défilés militaires". Le moteur effectue ensuite le lien entre la formulation de l'utilisateur en langage naturel, et les mots-clés indexés sur les séquences. Les résultats de recherche pourront par exemple présenter les séquences contenant les mots-clés suivants : "foule", "drapeau", "armée", "police". Une fois les résultats de recherche présentés, l'utilisateur peut cliquer sur une notice et accéder directement au média lié à une séquence spécifique, au sein d'une grille de recherche. Cela permet une large découvrabilité des archives, concrétisant ainsi les besoins utilisateurs fondateurs du modèle FRBR, de "trouver, identifier, sélectionner, obtenir" \footcite{zotero-item-653}. 
Les résultats de recherche pourront présenter en vignette, l’image clé associée au segment décrit, sur laquelle l’utilisateur peut cliquer. 
Outre la recherche textuelle, l'indexation fine des contenus audiovisuels peut développer l'usage d'autres modes de recherche visuels innovants. Le moteur Snoop, développé conjointement par l'\gls{ina} et l'Institut national de recherche en informatique et en
automatique (\gls{inria}), permet de croiser une exploration fine des contenus audiovisuels avec une interactivité qui place l'utilisateur au centre \footcite{zotero-item-744}. Fondé sur un bouclage de pertinence, il permet, à partir d'une requête sur une caractéristique visuelle, de proposer à l'utilisateur des résultats pertinents que celui-ci peut annoter comme résultats pertinents ou non pertinents \footcite{zotero-item-796}. Le moteur prend ces annotations en compte afin d'affiner la pertinence des résultats de recherche.
Ce moteur innovant suit la même logique que l'interface de validation pensée sur le modèle d'Amalia.js, et que le mode de recherche en \gls{nlp}, qui prévoit d'être affiné en fonction des retours utilisateurs. Les retours utilisateur, intégrés à l'interface, ne srevent pas seulement à valider les données mais améliorent le fonctionnement du logiciel, fond sur des processus automatisés. 

Conclusion sous partie 
\textit{In fine}, l'intégration de l'indexation automtique au sein d'un logiciel de gestion des collections exige une réadaptation du modèle de données via la création d'une nouvelle entité qui permet d'accueillir adéquatement les nouvelles métadonnées produites. De plus, l'implémentation du processus automatique implique la mise en place d'une interface de validation qui place l'utilisateur au centre de l'indexation définitive. Une fois indexées, la masse de données produites permet de concrétiser la finalité du FRBR, celle de la découvrabilité, via des outils de recherche innovants qui permettent d'explorer directement le contenu de l'image. 